% !TeX root = ../main.tex
\section{Прямые произведения, центр и автоморфизмы}

\task{60.4}{Когда $\Z_m\oplus\Z_n$ циклическая?}
\utv{}{Для положительных $m,n$ группа $\Z_m\oplus\Z_n$
 циклическая тогда и только тогда, когда $\gcd(m,n)=1$.}{
 Для любого $(a,b)$ его порядок делит $\lcm(m,n)$,
 поскольку эта степень одновременно обнуляет обе координаты.
 Если $\gcd(m,n)>1$, то
 $\lcm(m,n)=mn/\gcd(m,n)<mn=|\Z_m\oplus\Z_n|$.
 Элемента порядка всей группы нет, поэтому группа не циклическая.
 Если $\gcd(m,n)=1$, элемент $(\overline1,\overline1)$
 имеет порядок $\lcm(m,n)=mn$ и порождает всю группу.
}

\task{60.6}{Разложение $\C^*$}
\utv{}{Как группы с умножением, $\C^*\cong\Rp\times U$.}{
 Зададим $\varphi(z)=(|z|,z/|z|)$.
 Мультипликативность модуля даёт
 \[\varphi(zw)=\left(|z||w|,\frac z{|z|}\frac w{|w|}\right)
 =\varphi(z)\varphi(w).\]
 Обратное отображение --- $(r,u)\mapsto ru$.
 Действительно, $|ru|=r$ при $r>0$, $|u|=1$.
 Поэтому $\varphi$ --- изоморфизм.
}
\editnote{Аргумент комплексного числа определён по модулю $2\pi$.
 Интервал $[0,2\pi)$ с обычным сложением не является группой;
 корректное альтернативное описание ---
 $U\cong\R/(2\pi\Z)$ посредством $\theta\mapsto e^{i\theta}$.}

\task{60.15}{Центр прямого произведения}
\utv{Пункт а}{$Z(A\times B)=Z(A)\times Z(B)$.}{
 Условие, что $(a,b)$ коммутирует с каждым $(x,y)$,
 записывается как $(ax,by)=(xa,yb)$.
 Оно равносильно $ax=xa$ для всех $x\in A$ и
 $by=yb$ для всех $y\in B$, то есть
 $a\in Z(A)$, $b\in Z(B)$.
}
\utv{Пункт б}{Если $N\trianglelefteq A\times B$ и
 $N\cap(A\times\{e\})=N\cap(\{e\}\times B)=\{(e,e)\}$,
 то $N\le Z(A\times B)$.}{
 Пусть $(a,b)\in N$, $x\in A$.
 По нормальности $(x,e)(a,b)(x^{-1},e)\in N$.
 Умножим этот элемент на $(a,b)^{-1}\in N$:
 \[(xax^{-1}a^{-1},e)\in N\cap(A\times\{e\}).\]
 По условию он равен $(e,e)$, значит $xa=ax$ для всех $x$.
 Аналогично сопряжение элементом $(e,y)$ показывает $yb=by$
 для всех $y\in B$. Поэтому $(a,b)\in Z(A)\times Z(B)$,
 а по пункту а это центр произведения.
}
\editnote{Само сопряжение $(xax^{-1},b)$ не лежит в первом
 прямом сомножителе. Для применения условия о пересечении
 необходимо дополнительно умножить на $(a,b)^{-1}$.}

\task{60.17}{Прямое слагаемое и идемпотентная проекция}
Группа $G$ здесь абелева; используем аддитивную запись.
\utv{}{Подгруппа $A\le G$ является прямым слагаемым
 тогда и только тогда, когда существует сюръективный гомоморфизм
 $\pi:G\to A$ с $\pi^2=\pi$.
 В композиции $\pi^2$ используется включение $A\hookrightarrow G$.}{
 Если $G=A\oplus B$, каждый $g$ единственным образом представим
 как $g=a+b$. Положим $\pi(a+b)=a$.
 Единственность делает отображение корректным;
 покоординатное сложение доказывает, что это гомоморфизм.
 Оно сюръективно и фиксирует $A$, поэтому $\pi^2=\pi$.

 Обратно, пусть такая $\pi$ существует. Для любого $a\in A$
 по сюръективности найдётся $g$ с $\pi(g)=a$, откуда
 $\pi(a)=\pi^2(g)=\pi(g)=a$.
 Положим $B=\ker\pi$. Для любого $g\in G$
 \[g=\pi(g)+(g-\pi(g)),\qquad
 \pi(g-\pi(g))=\pi(g)-\pi^2(g)=0.
 \]
 Значит $G=A+B$. При $x\in A\cap B$ одновременно
 $\pi(x)=x$ и $\pi(x)=0$, так что $x=0$.
 Следовательно, $G=A\oplus B$.
}

\task{58.24}{Центры матричных групп}
\textbf{а) $\GL_n(\R)$.}
\utv{}{$Z(\GL_n(\R))=\{\lambda I:\lambda\in\R^*\}$.}{
 Скалярные матрицы коммутируют со всеми матрицами.
 Пусть $A=(a_{ij})$ центральна. Для любой обратимой диагональной
 матрицы $D=\diag(d_1,\dots,d_n)$ равенство $AD=DA$ даёт
 \[(d_j-d_i)a_{ij}=0.\]
 Для $i\ne j$ можно выбрать $d_i\ne d_j$, поэтому $a_{ij}=0$.
 Значит $A=\diag(a_1,\dots,a_n)$.
 При $i\ne j$ матрица $I+E_{ij}$ обратима, поскольку
 $E_{ij}^2=0$ и её обратная --- $I-E_{ij}$.
 Из $A(I+E_{ij})=(I+E_{ij})A$ следует
 $(a_i-a_j)E_{ij}=0$, то есть $a_i=a_j$.
 Все диагональные элементы равны ненулевому $\lambda$.
 При $n=1$ результат следует непосредственно.
}

\textbf{ж) Группа $T_n(\R)$ обратимых верхнетреугольных матриц.}
Все диагональные матрицы $D$ и матрицы $I+E_{ij}$ при $i<j$
принадлежат $T_n(\R)$. Поэтому то же доказательство сначала
обнуляет все внедиагональные элементы центральной матрицы,
а затем уравнивает её диагональные элементы. Следовательно,
\[Z(T_n(\R))=\{\lambda I:\lambda\in\R^*\}.\]
\editnote{Ответ $\lambda I+\mu E_{1n}$ неверен для этой группы
 при $n\ge2$ и $\mu\ne0$:
 для $D=\diag(2,1,\dots,1)$ имеем
 $DE_{1n}=2E_{1n}$, но $E_{1n}D=E_{1n}$.
 Исходное условие не уточняет поле в пункте ж; здесь, как в пункте а,
 выбрано $\R$. Доказательство работает над любым полем с более чем
 двумя элементами. Над $\F_2$ все обратимые диагонали единичны,
 группа становится унитреугольной, и вывод требует отдельного рассмотрения.}

\task{58.21}{Центральность $(ab)^2$}
\utv{}{Если $G=\langle a,b\rangle$,
 $a^2=b^2=(ab)^4=e$, то $(ab)^2\in Z(G)$.}{
 Положим $r=ab$. Так как $a^{-1}=a$, $b^{-1}=b$, имеем
 \[ara^{-1}=a(ab)a=ba=r^{-1},\qquad
 brb^{-1}=b(ab)b=ba=r^{-1}.
 \]
 Поэтому оба образующих переводят $r^2$ сопряжением в $r^{-2}$.
 Из $r^4=e$ следует $r^{-2}=r^2$.
 Значит $r^2$ коммутирует с $a,b$, а потому с любым словом
 в образующих и их обратных, то есть со всей $G$.
}
\editnote{Из $a^2=b^2=e$ не следует $ab=ba$.
 В исходнике это равенство использовалось без основания.}

\task{55.23}{Автоморфизмы аддитивной группы $\Q$}
\utv{}{Все автоморфизмы $(\Q,+)$ имеют вид
 $\varphi_a(x)=ax$, $a\in\Q^*$. Следовательно,
 $\Aut(\Q,+)\cong(\Q^*,\cdot)$.}{
 При $a\ne0$ отображение аддитивно и имеет обратное $x\mapsto x/a$,
 так что это автоморфизм.
 Пусть теперь $\varphi$ --- любой автоморфизм и $a=\varphi(1)$.
 Аддитивность даёт $\varphi(m)=ma$ для всех $m\in\Z$,
 включая отрицательные $m$.
 Для $m\in\Z$, $n\ge1$ имеем
 \[n\varphi(m/n)=\varphi(m)=ma,
 \qquad\varphi(m/n)=a\frac mn.\]
 Это определяет $\varphi$ на всей $\Q$.
 Если $a=0$, отображение нулевое и не биективно.
 Наконец, композиция $\varphi_a\circ\varphi_b=\varphi_{ab}$,
 что доказывает групповой изоморфизм.
}

\task{55.32, б}{Автоморфизмы $\Z_p$}
Пусть $p$ простое. Гомоморфизм $\Z_p\to\Z_p$ задаётся
образом $\overline1$, поэтому имеет вид $\varphi_a(x)=ax$.
При $a\ne0$ класс $a$ обратим, и обратное отображение
$x\mapsto a^{-1}x$ показывает биективность.
При $a=0$ отображение не является автоморфизмом.
Таким образом,
\[\Aut(\Z_p)\cong\Z_p^*,\qquad|\Aut(\Z_p)|=p-1.\]
Изоморфизм сопоставляет $a$ умножению на $a$;
сохранение операции следует из $\varphi_a\varphi_b=\varphi_{ab}$.
Пункт е этой задачи о группе $\Q$ уже решён в 55.23.

\task{57.40}{Автоморфизмы $S_3$ и $V_4$}
\textbf{а) $S_3$.}
Автоморфизм сохраняет порядки элементов, значит переставляет
три транспозиции. Действие на них задаёт инъективный
гомоморфизм $\Aut(S_3)\hookrightarrow S_3$:
если автоморфизм фиксирует транспозиции, он фиксирует их произведения,
а транспозиции порождают $S_3$. Поэтому $|\Aut(S_3)|\le6$.
С другой стороны, внутренние автоморфизмы задаются сопряжениями.
Ядро гомоморфизма
\[S_3\to\Inn(S_3),\qquad g\mapsto(x\mapsto gxg^{-1})\]
равно $Z(S_3)$. Центр тривиален:
транспозиция не коммутирует с другой транспозицией,
а тройной цикл не коммутирует с транспозицией,
которая переводит его в обратный цикл.
Следовательно, $|\Inn(S_3)|=6$.
Поскольку $\Inn(S_3)\le\Aut(S_3)$, получаем
\[\Aut(S_3)=\Inn(S_3)\cong S_3.\]

\textbf{б) $V_4$.}
Пусть $V_4=\{e,a,b,c\}$, где
$a^2=b^2=c^2=e$, $ab=c$, $bc=a$, $ca=b$.
Автоморфизм фиксирует $e$ и переставляет $a,b,c$.
Любая их перестановка сохраняет таблицу умножения:
произведение одинаковых неединичных элементов равно $e$,
а двух различных --- третьему. Поэтому каждая из $3!$ перестановок
задаёт автоморфизм, и $\Aut(V_4)\cong S_3$.
Поскольку $V_4$ абелева, $gxg^{-1}=x$ для всех $g,x$.
Значит $\Inn(V_4)=\{\operatorname{id}\}$.

\task{57.45}{Автоморфизмы диэдральной группы $D_4$}
Используем представление
\[D_4=\langle r,s\mid r^4=s^2=e,\ srs=r^{-1}\rangle.\]
Только $r,r^3$ имеют порядок $4$, поэтому автоморфизм переводит
$r$ в $r^u$, где $u\in\{1,3\}$.
Подгруппа $\langle r\rangle$ при этом сохраняется;
образ $s$ должен лежать вне неё, иначе образы не породят $D_4$.
Следовательно, $s\mapsto r^v s$, $v\in\Z_4$.
Каждое отображение
\[f_{u,v}(r)=r^u,\qquad f_{u,v}(s)=r^v s\]
сохраняет соотношения: $(r^v s)^2=e$ и
$(r^v s)r^u(r^v s)^{-1}=r^{-u}$.
Его образы порождают $D_4$, поэтому оно сюръективно;
для конечной группы это означает биективность.
Итак, имеется ровно $8$ автоморфизмов. Правило композиции:
\[f_{u,v}\circ f_{u',v'}=f_{uu',\,v+uv'},\]
где обе координаты считаются по модулю $4$.

Положим $\alpha=f_{1,1}$, $\beta=f_{-1,0}$.
Тогда $\alpha^4=\beta^2=e$, $\beta\alpha\beta=\alpha^{-1}$.
Четыре $\alpha^v=f_{1,v}$ и четыре $\alpha^v\beta=f_{-1,v}$
исчерпывают все автоморфизмы.
Это ровно диэдральные соотношения и нормальные формы, поэтому
\[\Aut(D_4)\cong D_4.\]

Найдём внутренние автоморфизмы. Сопряжение элементами $r,s$ даёт
\[\operatorname{conj}_r=f_{1,2}=\alpha^2,\qquad
 \operatorname{conj}_s=f_{-1,0}=\beta.\]
Поскольку $r,s$ порождают $D_4$, все внутренние автоморфизмы
порождаются этими двумя. Значит
\[\Inn(D_4)=\langle\alpha^2,\beta\rangle
 =\{f_{1,0},f_{1,2},f_{-1,0},f_{-1,2}\}\cong V_4.\]
Также $Z(D_4)=\{e,r^2\}$: вращение $r^j$ коммутирует с $s$
ровно при $2j\equiv0\pmod4$, а отражение не коммутирует с $r$.
Поэтому $\Inn(D_4)\cong D_4/Z(D_4)$ имеет порядок $4$.
\editnote{В представлении $D_4$ должно стоять $r^4=e$,
 а не $r^2=e$. Также для автоморфизма $s\mapsto r s$
 его квадрат переводит $s$ в $r^2s$, а не в $s$.}
